\documentclass[11pt]{article}
\usepackage{amsmath, amssymb}
\usepackage{geometry}
\usepackage{booktabs}
\usepackage{array}
\usepackage{parskip}
\usepackage{natbib}
\geometry{margin=2.5cm}

\title{\textbf{Novel GNN Framework: Four MLPs\\for Brain-Transformer-Aligned Neural Dynamics}}
\author{NeuralGraph Project}
\date{\today}

\begin{document}
\maketitle

%-----------------------------------------------------------------------
\section{Starting Point: Current GNN}

The current framework learns neural dynamics as:
\begin{equation}
\tau_i\frac{dv_i(t)}{dt} = -v_i(t) + V_i^{\text{rest}}
+ \sum_{j\in\mathcal{N}_i} W_{ij}\,
  \mathrm{ReLU}\!\big(v_j(t)\big)
+ I_i(t)
\label{eq:sim_current}
\end{equation}

for simulation, and
\begin{equation}
\frac{\widehat{dv}_i(t)}{dt}
= f_\theta\!\Big(
v_i(t),\,\mathbf{a}_i,\,
\sum_{j\in\mathcal{N}_i} \hat{W}_{ij}\,
g_\phi\!\big(v_j(t),\mathbf{a}_j\big)^2,\,
I_i(t)\Big)
\label{eq:gnn_current}
\end{equation}

for the GNN. The critical limitation is that both equations depend only on the
\emph{current} scalar value $v_j(t)$ --- no history, no temporal structure, no spatial structure,
no dynamic modulation of the receiver.
The two missing structures are not independent: the temporal delay between
two neurons is simply proportional to their physical distance, so adding
spatial positions automatically determines the relevant timescale of past
activity to consider.

%-----------------------------------------------------------------------
\section{History Vectors}

Each neuron now carries a history vector of its last $T$ voltage values:
\begin{equation}
\mathbf{v}_i(t) = \bigl[v_i(t),\; v_i(t-1),\; \ldots,\; v_i(t-T+1)\bigr]^\top
\in \mathbb{R}^T
\end{equation}

The scalar $v_i(t) = \mathbf{v}_i(t)[0]$ is the first element, so
the current framework is the special case $T=1$. The history vector
$\mathbf{v}_i \in \mathbb{R}^T$ plays the same role as the token embedding
$\mathbf{x}_i \in \mathbb{R}^d$ in a transformer \citep{vaswani2017}: it is
the raw representation from which all projections are computed, with $T$
playing the role of $d$.

%-----------------------------------------------------------------------
\section{Updated Equations}

\subsection*{Simulation}

The ground-truth simulation uses a per-edge axonal delay kernel
$G_{ij} \in \mathbb{R}^T$ that weights past activity according to the
travel time between neurons, and a per-neuron adaptation kernel
$Q_i \in \mathbb{R}^T$ that suppresses the response when the neuron
has been over-activated:

\begin{equation}
\boxed{
\tau_i\frac{dv_i(t)}{dt} = -v_i(t) + V_i^{\text{rest}}
+ \sigma\!\left(Q_i \cdot \mathbf{v}_i(t)\right)
\cdot \sum_{j\in\mathcal{N}_i} W_{ij}\,
  \mathrm{ReLU}\!\left(G_{ij} \cdot \mathbf{v}_j(t)\right)
+ I_i(t)
}
\label{eq:sim_new}
\end{equation}

In a special case, the dot products are over the $T$ time lags and
\begin{equation}
G_{ij}(k) = \exp\!\left(-\frac{\bigl(k - d_{ij}\bigr)^2}{2\sigma_g^2}\right),
\qquad d_{ij} = \frac{\|r_i - r_j\|}{v\,\Delta t}
\end{equation}

is a Gaussian bump peaked at the axonal travel time expressed in timesteps,
with neuron positions $r_i, r_j \in \mathbb{R}^3$, conduction velocity $v$,
and timestep $\Delta t$. The index $k=0$ corresponds to the current time
$v_j(t)$ and $k=T-1$ to the furthest past $v_j(t-T+1)$.
Note that space and time are not independent here: the delay $d_{ij}$
is directly proportional to the physical distance $\|r_i - r_j\|$,
so knowing neuron positions automatically determines how far back in
the history $\mathbf{v}_j$ the signal from $j$ arrives at $i$.
The dot product $G_{ij} \cdot \mathbf{v}_j$ is therefore a special case
of a receiver-dependent temporal filter: the kernel shape is fully
determined by the distance between $i$ and $j$, which is itself encoded
in the identity of the receiver $i$.

\subsection*{GNN}

The GNN separates the message computation into two distinct and
independently learned stages. The first stage compresses the sender's
history $\mathbf{v}_j \in \mathbb{R}^T$ into a scalar $\tilde{v}_j \in \mathbb{R}$
through a temporal filter $h_\xi$ that depends only on the receiver
identity $\mathbf{a}_i$ and the sender history $\mathbf{v}_j$.
The receiver embedding $\mathbf{a}_i$ implicitly encodes spatial context
including the receiver's position, so distance-dependent delay is learned
without being explicitly imposed. The second stage applies a nonlinear
message function $g_\phi$ to the filtered scalar, conditioned on both
neuron embeddings. This separation is the key architectural novelty:
$h_\xi$ is a pure spatio-temporal filter; $g_\phi$ is a pure nonlinearity
with heterogeneity. They are independent networks with independent
parameter sets $\xi$ and $\phi$.

Defining:
\begin{equation}
\tilde{v}_j(t) = h_\xi\!\left(\mathbf{a}_i,\,\mathbf{v}_j(t)\right) \in \mathbb{R}
\end{equation}

the full GNN dynamics are:

\begin{equation}
\boxed{
\frac{\widehat{dv}_i(t)}{dt}
= f_\theta\!\left(
\mathbf{v}_i(t),\;\mathbf{a}_i,\;
q_\psi\!\left(\mathbf{a}_i,\,\mathbf{v}_i(t)\right)
\cdot \sum_{j\in\mathcal{N}_i}
\hat{W}_{ij}\;
g_\phi\!\left(\mathbf{a}_i,\,\mathbf{a}_j,\,\tilde{v}_j(t)\right),\;
I_i(t)
\right)
}
\label{eq:gnn_new}
\end{equation}

The four MLPs have fully independent parameter sets and clearly separated
roles:

\begin{center}
\begin{tabular}{lllp{5.8cm}}
\toprule
MLP & Params & Input dim & Role \\
\midrule
$h_\xi$ & $\xi$ & $d_e + T$ &
  Temporal filter: receiver-dependent compression of sender history
  into a scalar. Generalises $G_{ij} \cdot \mathbf{v}_j$.
  Distance encoded implicitly through $\mathbf{a}_i$. \\
$g_\phi$ & $\phi$ & $2d_e + 1$ &
  Nonlinear message: applies learned nonlinearity and heterogeneity
  to the filtered scalar $\tilde{v}_j$. Generalises $\mathrm{ReLU}(\cdot)$. \\
$q_\psi$ & $\psi$ & $d_e + T$ &
  Adaptation gate $\in(0,1)$: neuron-dependent suppression
  of incoming drive via receiver history and identity. \\
$f_\theta$ & $\theta$ & $T + d_e + 2$ &
  Node update: integrates history, identity, aggregated
  message, and external drive. \\
\bottomrule
\end{tabular}
\end{center}

The separation of $h_\xi$ from $g_\phi$ allows $h_\xi$ to specialise entirely
in temporal pattern extraction and $g_\phi$ to specialise in nonlinear synaptic
transformation, each with its own inductive bias.

%-----------------------------------------------------------------------
\section{Relationship to Transformers}

\subsection*{The transformer Q-K-V structure}

The scaled dot-product attention mechanism is \citep{vaswani2017}:
\begin{equation}
s_{ij} = \frac{(\mathbf{x}_i W_Q)(\mathbf{x}_j W_K)^\top}{\sqrt{d_k}},
\qquad
\mathrm{output}_i = \sum_j \mathrm{softmax}_j(s_{ij})\cdot \mathbf{x}_j W_V
\label{eq:transformer}
\end{equation}

where $\mathbf{x}_i \in \mathbb{R}^d$ is the token embedding and $W_Q$,
$W_K$, $W_V$ are learned projection matrices.
These are named Query, Key, and Value, reflecting their canonical roles:
$W_Q$ projects token $i$ into what it is searching for, $W_K$ projects
token $j$ into what it advertises, and $W_V$ projects token $j$ into what
it actually sends. The dot product $q_i \cdot k_j$ measures how well the
receiver's expectation matches the sender's content, and softmax normalises
this score into a routing weight over all senders. Here we map these roles
onto neural circuit mechanisms, which clarifies both the analogy and where
it must be abandoned in favour of biologically grounded MLPs.

Note that $W_Q$ and $W_K$ are mathematically interchangeable in
Eq.~\eqref{eq:transformer}: swapping them and transposing the score matrix
$S \to S^\top$ leaves the attention output unchanged. Their real purpose is
not role separation but breaking symmetry: a single matrix $W$ would give
$s_{ij} = s_{ji}$ (undirected attention), whereas two different matrices
produce asymmetric, directed attention. In the GNN, directionality is
already handled by $\hat{W}_{ij} \neq \hat{W}_{ji}$, making the $W_Q
\cdot W_K$ asymmetry trick entirely redundant.

\subsection*{Cortical anatomy as a transformer: the K\"{o}nig--Negrello mapping}

A complementary perspective on the transformer--neuroscience connection is
provided by \citet{konig2026}, who propose that the mammalian cerebral cortex
implements a transformer-like architecture through the laminar structure of
cortical columns. In their mapping, one cortical column plays the role of one
transformer token, and the laminar microcircuit implements the Q-K-V
operations through distinct anatomical pathways. The correspondence is
summarised in Table~\ref{tab:konig}.

\begin{table}[h]
\centering
\begin{tabular}{lll}
\toprule
Transformer concept & Cortical anatomy \citep{konig2026} & GNN analog \\
\midrule
Token & Cortical column & Neuron $i$ with embedding $\mathbf{a}_i$ \\
Input embedding & Thalamic drive to L4 & External input $I_i(t)$ \\
Values (V) & L4$\to$L2/3 feedforward & $\tilde{v}_j = h_\xi(\mathbf{a}_i, \mathbf{v}_j)$ \\
Keys (K) & L2/3 tangential connections & absorbed into $\hat{W}_{ij}$, dropped \\
Queries (Q) & L1 top-down apical feedback & $q_\psi(\mathbf{a}_i, \mathbf{v}_i)$ \\
Attention Q$\cdot$K & Dendritic coincidence detection & $q_\psi \cdot g_\phi$ multiplicative gate \\
Multi-head & Parallel cell subpopulations & $H$ copies of $h_\xi$, $g_\phi$, $q_\psi$ \\
FFN sublayer & Local interlaminar recoding & $f_\theta(\mathbf{v}_i, \mathbf{a}_i, \cdot)$ \\
Residual & Ongoing intracortical activity & $-v_i + V_i^{\text{rest}}$ leak term \\
LayerNorm & E/I gain control & implicit in $\sigma(\cdot)$ of $q_\psi$ \\
Self-attn connectivity & Horizontal + feedback coupling & $\hat{W}_{ij}$ sparse anatomy \\
\bottomrule
\end{tabular}
\caption{Mapping between transformer components, cortical anatomy
\citep{konig2026}, and the present GNN framework.}
\label{tab:konig}
\end{table}

The K\"{o}nig--Negrello mapping resolves the Q/K asymmetry question
raised above: in cortex, Q and K are \emph{anatomically separated}. Queries
arrive via L1 top-down feedback onto apical dendrites of L2/3 and L5
pyramidal neurons, while Keys are conveyed laterally through L2/3 tangential
horizontal projections. Because these are distinct anatomical pathways,
they are not interchangeable, in contrast to the abstract transformer
equation \citep{konig2026}. In the single-neuron GNN, this laminar
distinction has no substrate --- all inputs arrive at the same node --- which
further justifies absorbing K into the connectivity $\hat{W}_{ij}$ and
retaining only Q as the adaptation gate $q_\psi$.

The cortical interpretation of multi-head attention is also more specific
than the AMPA/NMDA analogy: \citet{konig2026} propose that distinct
frequency bands (gamma for encoding in L2/3, alpha/beta for decoding in L5)
implement different heads in parallel, with oscillatory synchrony providing
the dynamic gating. This is consistent with our $H$-head formulation but
suggests a richer temporal structure: heads parameterised by oscillation
frequency rather than receptor timescale alone. In early sensory cortex,
the authors further propose ``attention-free'' transformer variants where
routing is hardwired into horizontal connectivity --- exactly the regime
where $\hat{W}_{ij}$ is fixed and only $h_\xi$, $g_\phi$, $q_\psi$ are
learned.

\subsection*{Linear intermediate: three matrices}

Replacing the token embedding $\mathbf{x}_i \in \mathbb{R}^d$ with the
neuron history $\mathbf{v}_i \in \mathbb{R}^T$, the transformer Q-K-V
structure maps directly onto neural circuit mechanisms through three
learned matrices acting on the history vectors:

\begin{center}
\begin{tabular}{lllll}
\toprule
Matrix & Shape & Input & Output & Neural role \\
\midrule
$\hat{W}_Q$ & $\mathbb{R}^{T \times d_k}$ & $\mathbf{v}_i$ history of receiver
  & $q_i \in \mathbb{R}^{d_k}$ & what pattern does $i$ expect? \\
$\hat{W}_K$ & $\mathbb{R}^{T \times d_k}$ & $\mathbf{v}_j$ history of sender
  & $k_j \in \mathbb{R}^{d_k}$ & what pattern does $j$ contain? \\
$\hat{W}_V$ & $\mathbb{R}^{T \times 1}$   & $\mathbf{v}_j$ history of sender
  & $\tilde{v}_j \in \mathbb{R}$  & what does $j$ actually transmit? \\
\bottomrule
\end{tabular}
\end{center}

\begin{equation}
q_i = \mathbf{v}_i\,\hat{W}_Q, \quad
k_j = \mathbf{v}_j\,\hat{W}_K, \quad
\tilde{v}_j = \mathbf{v}_j\,\hat{W}_V, \quad
s_{ij} = \frac{q_i \cdot k_j}{\sqrt{d_k}}, \quad
m_{ij} = \hat{W}_{ij} \cdot \sigma(s_{ij}) \cdot \tilde{v}_j
\end{equation}

$\hat{W}_Q$ detects what temporal pattern the postsynaptic neuron is primed
to receive, analogous to the apical-dendritic Query signal in
\citet{konig2026}.
$\hat{W}_K$ detects what pattern the presynaptic neuron is currently
broadcasting, analogous to the L2/3 tangential Key pathway.
$\hat{W}_V$ extracts the scalar signal actually transmitted along the axon,
analogous to the L4$\to$L2/3 Value stream.
However this intermediate is \emph{more transformer-driven than biologically
driven} for three reasons.
All three projections are linear and cannot detect nonlinear temporal
patterns such as bursts or oscillations --- precisely the nonlinear
dendritic coincidence detection that \citet{konig2026} identify as the
biological substrate of the Q$\cdot$K dot product.
$\hat{W}_Q$ and $\hat{W}_V$ are shared across all neurons, making adaptation
and filtering neuron-independent in contradiction to known biological
heterogeneity \citep{konig2026}.
Most importantly, $\hat{W}_K$ has no clear single-neuron counterpart:
in cortex, the Key pathway is a population-level tangential connectivity
pattern, not a property of individual neurons, making it natural to absorb
into $\hat{W}_{ij}$ at the GNN level.

\subsection*{From linear matrices to nonlinear MLPs}

These limitations motivate replacing the three matrices with independent
nonlinear MLPs. $\hat{W}_V$ becomes $h_\xi(\mathbf{a}_i, \mathbf{v}_j)$:
a nonlinear receiver-dependent temporal filter that compresses the full
history into a scalar $\tilde{v}_j$, with $\mathbf{a}_i$ encoding the
receiver's spatial context and expected input timescale. The fixed
nonlinearity $\mathrm{ReLU}$ becomes $g_\phi(\mathbf{a}_i, \mathbf{a}_j,
\tilde{v}_j)$: a learned neuron-dependent transformation implementing the
dendritic nonlinearity of \citet{konig2026} in a data-driven form.
$\hat{W}_Q$ becomes $q_\psi(\mathbf{a}_i, \mathbf{v}_i)$: a nonlinear
adaptation gate neuron-dependent through $\mathbf{a}_i$, corresponding to
the top-down apical-dendritic Query signal whose strength depends on the
receiver's current state.
$\hat{W}_K$ is dropped entirely: the sparse anatomy $\hat{W}_{ij}$ already
encodes directed connectivity and replaces softmax routing, consistent with
the proposal of \citet{konig2026} that in early sensory cortex routing is
hardwired into horizontal connectivity rather than dynamically computed.

\subsection*{Multi-head extension}

The natural GNN analog of multi-head attention \citep{vaswani2017} runs
$H$ independent copies of $h_\xi$, $g_\phi$, and $q_\psi$:
\begin{equation}
\tilde{v}_j^{(h)} = h_\xi^{(h)}\!\left(\mathbf{a}_i,\,\mathbf{v}_j\right),
\qquad
m_{ij} = \hat{W}_{ij} \cdot \frac{1}{H}\sum_{h=1}^{H}
q_\psi^{(h)}\!\left(\mathbf{a}_i,\,\mathbf{v}_i\right) \cdot
g_\phi^{(h)}\!\left(\mathbf{a}_i,\,\mathbf{a}_j,\,\tilde{v}_j^{(h)}\right)
\end{equation}

At the synaptic level, the heads correspond to distinct receptor types
coexisting on the same synapse: AMPA ($\tau \sim 2\,\mathrm{ms}$),
NMDA ($\tau \sim 50\,\mathrm{ms}$), and GABA ($\tau \sim 10\,\mathrm{ms}$).
At the population level, \citet{konig2026} propose that distinct frequency
bands implement different heads: gamma oscillations ($\sim$40--80\,Hz) drive
the encoding (L2/3 Value) stream while alpha/beta ($\sim$8--28\,Hz) drives
the decoding (L5 Query) stream.
Both views predict multiple parallel temporal filters per synapse, which
the $H$-head formulation captures without committing to either mechanism.

\subsection*{Fundamental differences between the GNN and transformer}

Three properties of the GNN have no transformer analog.
The GNN solves a differential equation $\tau_i \dot{v}_i = \ldots$ with true
continuous time and persistent recurrent state; the transformer maps a
sequence to a sequence in a single feedforward pass with no dynamics.
\citet{konig2026} explicitly acknowledge this: their mapping is intended at
the level of effective input--output transformations, not as a claim about
the detailed temporal unfolding of cortical dynamics, and they identify
continuous recurrent dynamics as the principal open problem for their
framework. The GNN addresses this directly.
$\hat{W}_{ij}$ is sparse, anatomically interpretable, and permanently
learned; transformer attention weights are dense and recomputed for every
new input.
Each neuron carries a permanent latent identity $\mathbf{a}_i$ encoding
biophysical type across all timesteps and trials; transformer tokens have
no such persistent node state. This permanent embedding is the single-neuron
counterpart of the column-level identity implicit in the topographic maps
of \citet{konig2026}.

\subsection*{Summary table}

\begin{center}
\begin{tabular}{lll}
\toprule
Concept & Transformer \citep{vaswani2017} & GNN (this work) \\
\midrule
Sender filter & $\mathbf{v}_j \hat{W}_V$ linear &
  $h_\xi(\mathbf{a}_i, \mathbf{v}_j)$ nonlinear, receiver-dependent \\
Sender nonlinearity & fixed (none beyond $W_V$) &
  $g_\phi(\mathbf{a}_i, \mathbf{a}_j, \tilde{v}_j)$ learned \\
Receiver gate & $\mathbf{v}_i \hat{W}_Q$ linear, shared &
  $q_\psi(\mathbf{a}_i, \mathbf{v}_i)$ nonlinear, neuron-dependent \\
Sender matching & $\mathbf{v}_j \hat{W}_K$ linear &
  absorbed into $\hat{W}_{ij}$, dropped \\
Routing & softmax$_j(s_{ij})$ dense &
  $\hat{W}_{ij}$ sparse anatomy \\
Normalisation & $\sum_j a_{ij} = 1$ & unconstrained \\
Multi-head & $H$ linear Q-K-V projections &
  $H$ copies of $h_\xi$, $g_\phi$, $q_\psi$ \\
Input representation & $\mathbf{x}_i \in \mathbb{R}^d$ token &
  $\mathbf{v}_i \in \mathbb{R}^T$ history \\
Node identity & none & $\mathbf{a}_i$ permanent embedding \\
Time & positional encoding, static & continuous ODE, recurrent \\
Geometry & none & implicit in $\mathbf{a}_i$, explicit in $d_{ij}$ \\
\bottomrule
\end{tabular}
\end{center}

%-----------------------------------------------------------------------
\section{Summary}

The novel framework introduces four independent MLPs: $h_\xi$, $g_\phi$,
$q_\psi$, and $f_\theta$, alongside the connectivity matrix $\hat{W}_{ij}$.
The central architectural contribution is the separation of the message
computation into a temporal filter $h_\xi(\mathbf{a}_i, \mathbf{v}_j)$
and a nonlinear transformation $g_\phi(\mathbf{a}_i, \mathbf{a}_j,
\tilde{v}_j)$. The filter operates on $\mathbb{R}^T$ and specialises in
selecting which past activity of the sender matters to this receiver; the
nonlinearity operates on $\mathbb{R}$ and specialises in synaptic
transformation.

The framework is grounded in two complementary perspectives. From the
bottom up, each MLP has a direct biological counterpart: $h_\xi$ generalises
the axonal delay kernel, $g_\phi$ generalises the synaptic nonlinearity,
$q_\psi$ generalises neural adaptation, and $f_\theta$ generalises the
membrane integration. From the top down, the framework maps cleanly onto
the transformer Q-K-V structure as nonlinearised and anatomically separated
versions of the linear projections, consistent with the cortical column
interpretation of \citet{konig2026}: $q_\psi$ is the apical-dendritic
Query, $h_\xi$ is the Value stream, $g_\phi$ is the dendritic nonlinearity,
and $\hat{W}_{ij}$ absorbs the Key pathway by encoding directed connectivity
explicitly. The principal advance over \citet{konig2026} is that the present
framework operates at the level of individual neurons with continuous
dynamics, learned connectivity, and spatial geometry --- the three elements
their mapping explicitly defers to future work.

%-----------------------------------------------------------------------
\bibliographystyle{plainnat}
\begin{thebibliography}{9}

\bibitem[K\"{o}nig and Negrello(2026)]{konig2026}
P.~K\"{o}nig and M.~Negrello.
\newblock The neuroscience of transformers.
\newblock arXiv:2603.15339 [q-bio.NC], 2026.

\bibitem[Vaswani et~al.(2017)]{vaswani2017}
A.~Vaswani, N.~Shazeer, N.~Parmar, J.~Uszkoreit, L.~Jones, A.~N. Gomez,
  \L.~Kaiser, and I.~Polosukhin.
\newblock Attention is all you need.
\newblock In \emph{Advances in Neural Information Processing Systems},
  volume~30, 2017.

\end{thebibliography}

\end{document}
